\section{Alexandrov topology and the Husainov comparison}

\begin{theorem}[The Alexandrov topology recovers the poset]\label{thm:alexandrov-recovers-poset}\lean{Obsmetrologia.Alexandrov.alexandrov_specialization_recovers_poset}\leanok\uses{thm:kle}
For the Alexandrov topology of a poset $(P,\le)$ --- opens are the up-sets --- the opens are
closed under \emph{arbitrary} intersection and union (the defining Alexandrov property), the
specialization order recovers $\le$ exactly ($x \rightsquigarrow y \iff x \le y$, as predicates
$\rightsquigarrow\ =\ \le$), and a map is continuous iff it is monotone. This is the classical
equivalence \emph{Alexandrov spaces $\cong$ preorders}, proved over the derived poset
infrastructure ($\texttt{alexandrov\_specialization\_recovers\_poset}$).
\end{theorem}
\begin{theorem}[\v{C}ech cohomology vanishes on the bounded sub-sites]\label{thm:husainov-cech-acyclic}\lean{Obsmetrologia.Alexandrov.husainov_cech_acyclic_on_bounded_sites}\leanok\uses{thm:h1-zero-iff-confluent, thm:alexandrov-recovers-poset}
On the bounded sub-sites --- the strict ($\kappa<\omega$) / Scott ($\kappa=\omega$) regime --- the
\v{C}ech $H^1$ of every cover vanishes ($\texttt{CechAcyclic}$, reusing
\cref{thm:h1-zero-iff-confluent}); no \v{C}ech complex above degree one is constructed. This is a
statement about the \v{C}ech computation \emph{alone}; it makes no comparison with the derived
functor. In this model the \v{C}ech vanishing is in fact unconditional --- every cover glues ---
so the regime here gates only the general dichotomy, where $H^1$ may be nonzero in the non-Scott
regimes.
\end{theorem}
\begin{theorem}[Husainov: \v{C}ech $=$ derived-functor cohomology]\label{thm:husainov-cech-derived}\notready\uses{thm:husainov-cech-acyclic, thm:alexandrov-recovers-poset}
\emph{(Open frontier --- stated, not formalized.)} On an Alexandrov poset, \v{C}ech cohomology
agrees with the \emph{derived-functor} sheaf cohomology of the sheafification,
$\check{H}^n(F_U) \cong R^n\Gamma(F_U)$, exactly when the poset's Dedekind--MacNeille cuts are
acyclic --- a necessary-and-sufficient condition, algorithmically computable for finite posets
([@husainov2026]). The load-bearing content is the comparison itself, and it requires
\emph{derived-functor} cohomology $R^n\Gamma$ --- the right-derived functors of global sections,
built from an injective or flasque resolution of $F_U$. That resolution machinery is not
available in the dependency-free core, and without it there is no honest non-trivial Lean
statement of the comparison: defining $R\Gamma := \check{H}$ makes it a tautology (the exact
defect this pass removed), while conditioning it on an unconstrained abstract $R\Gamma$ makes it
refutable. This node therefore stays \textbf{red} until the resolution machinery is formalized;
it is the genuine frontier this corner of the library reaches, recorded here at full strength
rather than faked in Lean. \cref{thm:husainov-cech-acyclic} supplies the \v{C}ech half it will
build on.
\end{theorem}

\begin{theorem}[Hole/diamond cuts isolate a Husainov special case]\label{thm:hole-diamond-husainov-frontier}\lean{Obsmetrologia.Alexandrov.HoleDiamondCutFamily.holeDiamond_husainov_comparison_of_cut_acyclic}\leanok\uses{thm:husainov-cech-acyclic, thm:alexandrov-recovers-poset}
For a finite listed causal set whose readout region is a single closed causal diamond, the
formal carrier $\texttt{HoleDiamondCutFamily}$ packages the concrete
$\texttt{HoleRegion}(C,R)$ / $\texttt{ReadoutExterior}(C,R)$ cut family. Its interior is
future-closed, its exterior is past-closed, the two sides are disjoint, and no interior event
causally precedes an exterior event. The Lean theorem here does not define derived-functor
cohomology or Dedekind--MacNeille cuts; instead it proves the exact frontier reduction:
for any future formal predicates $\texttt{CutAcyclic}$ and $\texttt{Comparison}$, a genuine
Husainov criterion
$\forall P,\ \texttt{CutAcyclic}(P) \to \texttt{Comparison}(P)$ turns cut acyclicity of this
hole/diamond family into the corresponding comparison statement for this family.
\end{theorem}

\begin{theorem}[The Husainov readout and the causal-net diamond region agree]\label{thm:closed-diamond-readout-eq-diamond-region}\lean{Obsmetrologia.Alexandrov.closedDiamondReadout_asProp_eq_diamondRegion}\leanok\uses{thm:hole-diamond-husainov-frontier, thm:causal-net}
The Boolean closed-diamond readout $\texttt{ClosedDiamondReadout}$ that
$\texttt{HoleDiamondCutFamily}$ builds its special Husainov family from, read as a Prop-valued
predicate, is pointwise equal to $\texttt{Causal.Net.DiamondRegion}$, the causal-net chapter's
named closed causal diamond region. The Husainov special case and the causal-net region theory
are two presentations --- Boolean-decidable and Prop-valued --- of the same closed causal
diamond, not two different constructions that happen to share a name.
\end{theorem}

\begin{theorem}[Hole-region causal diamonds carry a sequencing kinematic space]\label{thm:hole-diamond-kinematic-space}\lean{Obsmetrologia.CausalSets.BlackHoles.holeDiamond_kinematic_space}\leanok\uses{thm:hole-diamond-husainov-frontier}
The kinematic space of finite causal diamonds is the causal set whose events are
$\texttt{DiamondPoint}(C)$ pairs and whose order is endpoint sequencing:
$(a,b)\prec(a',b')$ exactly when $b\le a'$.  The tempting nesting relation is not this
causal set, because nesting has a self-edge at every diamond point and therefore cannot
carry the required strict-order closure.  Sequencing is transitive and irreflexive, so it
does form the derived causal set of diamonds.  Consequently any readout region on this
diamond space inherits the same no-readout/exterior cut laws as an ordinary finite causal
set: the hole region is future-closed, the readout exterior is past-closed, the sides are
disjoint, and no hole-region diamond causally precedes an exterior diamond.  This is the
proof-theoretic point at which a kinematic point, kinematic region, and kinematic patch can
be read without importing a background continuum: points are \(\texttt{DiamondPoint}\) values,
regions are Boolean predicates on the sequencing causal set \(K(C)\), and patches are
no-readout subregions carried by the same finite patch laws used downstream.
\end{theorem}

\begin{theorem}[Deshrouding: a shroud event stands as a diamond's upper endpoint]\label{thm:deshrouded-diamond-point}\lean{Obsmetrologia.CausalSets.BlackHoles.deshroudedDiamondPoint, Obsmetrologia.CausalSets.BlackHoles.deshroudedDiamondPoint_upper_not_exterior}\leanok\uses{thm:hole-diamond-kinematic-space}
A $\texttt{DiamondPoint}$ needs a genuine upper endpoint the lower endpoint strictly reaches.
Before the diamond's true readout is settled, a currently witnessed $\texttt{Shroud}$ event ---
inside the hole region, adjacent to the readout exterior, but proven not itself exterior --- can
stand in for it: $\texttt{deshroudedDiamondPoint}$ builds a genuine $\texttt{DiamondPoint}$ from
any such shroud event, and its upper endpoint is proven never to lie in the readout exterior, so
the proxy never overclaims completion it has not reached.
\end{theorem}

\begin{theorem}[Deshrouded diamonds only grow]\label{thm:deshrouded-diamond-point-nested}\lean{Obsmetrologia.CausalSets.BlackHoles.deshroudedDiamondPoint_nested_of_causalLe}\leanok\uses{thm:deshrouded-diamond-point}
Deshrouding a later-reached shroud event nests the earlier deshrouded diamond under
$\texttt{nestingPrec0}$: the lower endpoint is unchanged and the upper endpoint only advances.
A live process's proxy diamond therefore only ever grows as more of it is witnessed; it never
shrinks or moves its admission.
\end{theorem}

\begin{theorem}[Kugelblitz local hole-entry law]\label{thm:kugelblitz-local-hole-entry}\lean{Obsmetrologia.CausalSets.BlackHoles.ArticulationHoleRegionSwitchesAt, Obsmetrologia.CausalSets.BlackHoles.articulationKugelblitz_holeRegionSwitch, Obsmetrologia.CausalSets.BlackHoles.KugelblitzArticulationSite.tail_not_in_hole, Obsmetrologia.CausalSets.BlackHoles.KugelblitzArticulationSite.next_not_exterior, Obsmetrologia.CausalSets.BlackHoles.KugelblitzArticulationSite.holeRegionSwitch}\leanok\uses{thm:hole-diamond-kinematic-space}
The part of the growth frontier proved so far is the local articulation switch,
not the global strict-growth equivalence.  In the causal set generated by the
next articulation, a kugelblitz step proves that the tail endpoint is exterior
and therefore not in the computed hole region, while the paired next
articulation lies on the shroud, is in the computed hole region, and is not
exterior.  This is the finite formation carrier proved beneath the open
conjecture that strict region growth happens exactly at such articulation
sites.
\end{theorem}

\begin{theorem}[Kugelblitz growth conjecture]\label{thm:kugelblitz-growth-conjecture}\lean{Obsmetrologia.CausalSets.BlackHoles.kugelblitz_growth_conjecture}\notready\uses{thm:kugelblitz-local-hole-entry}
\emph{(Open articulation-site frontier.)}  A kugelblitz is an articulation site: a
tail articulation together with the next utterance whose generated causal diamond has lower
endpoint in the readout exterior and new upper endpoint on the shroud.  Thus the diamond
crosses the readout horizon in the only finite direction locally allowed by the hole laws:
exterior to interior.  The conjecture says that, for a fixed readout predicate on the
articulation chain, the computed hole region strictly grows from a tail articulation to the
next paired articulation exactly at those kugelblitz sites.  The Lean declaration records
this global strict-growth equivalence as a proposition, not as a proved theorem; what is
proved so far is the local articulation-site switch.
\end{theorem}

\begin{theorem}[Hole/diamond Dedekind--MacNeille acyclicity]\label{thm:hole-diamond-cut-acyclicity}\lean{Obsmetrologia.Alexandrov.holeDiamond_cut_acyclicity_conjecture}\notready\uses{thm:hole-diamond-husainov-frontier}
\emph{(Open special case.)} The intended special-case obligation is that the
Dedekind--MacNeille cuts attached to every $\texttt{HoleDiamondCutFamily}$ are acyclic. This is
the concrete tractable subproblem of \cref{thm:husainov-cech-derived}: proving it would make the
Husainov comparison available for exactly the finite hole/readout pairs whose readout side is a
single connected causal-diamond-generated region, without proving the general comparison node.
\end{theorem}

\begin{theorem}[Tarski sheaf closure identifies the harmonic kernel with $H^0$]\label{thm:tarski-sheaf-closure}\lean{Obsmetrologia.Tarski.harmonic_kernel_h0_iso}\leanok\uses{thm:closure-operator, thm:op-least, thm:settling-hodge-iso}
For a directed cell complex valued in a complete lattice, the Tarski drive has a
least fixed point by the closure/fixed-point spine. Its harmonic kernel and its
degree-zero settled-section object are equivalent. This is the shared closure
carrier used later by keratin contacts, liquid short-range order, and graph
screening.
\end{theorem}
